

Turpmākajā darbā iespējami vairāki virzieni piedāvāto \Gls{snn} modeļu un uz \Gls{fpga} balstītā \Gls{snn} paātrinātāja veiktspējas, mērogojamības, efektivitātes un piemērojamības uzlabošanai.

Pirmkārt, dažos uzdevumos impulsu modeļu precizitāte joprojām atpaliek no tradicionālajiem \Gls{ann}. Turpmāk varētu pētīt progresīvākas \Gls{snn} arhitektūras un apmācības metodes; paātrinātāja paplašināšana atlikuma savienojumu vai uz transformeriem balstītu modeļu atbalstam palielinātu tā piemērojamību dažādām \Gls{snn} arhitektūrām.

Otrkārt, varētu veikt kopīgu neironu tīkla un aparatūras arhitektūras optimizāciju. Turpmākajās realizācijās varētu pētīt neironu stāvokļu un starpposmu pazīmju karšu saspiešanu. Saspiestu neironu stāvokļu ar mazu bitu platumu glabāšana mikroshēmā varētu samazināt ar atmiņas piekļuvi saistīto jaudas patēriņu. Šim nolūkam varētu izmantot evolucionārās apmācības metodes un aparatūru ievērojošu \Gls{nas}, lai meklētu tīkla konfigurācijas ar labāku kompromisu starp precizitāti, resursu izmantojumu un enerģijas patēriņu.

Līdzīgi dažus no šiem ierobežojumiem var mazināt, integrējot lielāku datu apstrādes konveijera daļu \Gls{fpga} paātrinātājā vai ciešāk integrējot to ar attēlveidošanas sensoriem, samazinot dārgas pārsūtīšanas starp resursdatora procesoru, ārējo atmiņu un \Gls{fpga}. Tas varētu ietvert vokselizācijas, impulsu kodēšanas un objektu detektora galvas daļu aparatūras realizāciju. Tāpat LiDAR vai \Gls{spad} sensoru izeju tieša apstrāde \Gls{fpga} varētu samazināt starpposmu histogrammu veidošanu un \Gls{tdc} apstrādi, savukārt ciešāka integrācija ar notikumu kamerām varētu ļaut paātrinātājam tieši apstrādāt asinhronus notikumus.

Vēl viens turpmākais virziens ir sensoru sistēmas, \Gls{snn} modeļu un paātrinātāja integrācija automatizētas braukšanas vai citās iegultās skaitļošanas platformās. Lai gan pašreizējais darbs demonstrē reālu sensoru datu inferences iespējamību ar \Gls{snn} \Gls{fpga} aparatūrā, integrācija ar automatizētas braukšanas ietvariem ļautu novērtēt pilnu uztveres un lēmumu pieņemšanas konveijeru.

Pašreizējie rezultāti norāda uz lielu potenciālu turpmākajās \Gls{snn} realizācijās un to paātrināšanā ar \Gls{fpga}, īpaši tad, ja notikumu sensori integrēti pilnīgā sistēmā, kur uztvere, neirālā apstrāde un aparatūras arhitektūra projektētas kopīgi.
